\documentclass[10pt]{article}
\usepackage[letterpaper,hmargin=0.5in,vmargin=0.36in]{geometry}
\usepackage{longtable}
\usepackage{needspace}
\usepackage{multirow}
\usepackage{tikz-cd}
\usepackage{amsmath,amssymb,mathtools}
\usepackage{booktabs,array,colortbl}
\usepackage{xcolor}
\usepackage{pdfrender}
\usepackage[most]{tcolorbox}
\pagestyle{empty}
\setlength{\parindent}{0pt}
\definecolor{boxbg}{HTML}{F3F3F3}
\definecolor{boxfr}{HTML}{9A9A9A}
\definecolor{hdr}{HTML}{E4E4E4}
\newtcolorbox{sheetbox}{breakable,coltext=black,colback=boxbg,colframe=boxfr,boxrule=0.5pt,arc=2pt,left=6pt,right=6pt,top=2pt,bottom=2pt}
\newtcolorbox{smallbox}{breakable,fontupper=\small,coltext=black,colback=boxbg,colframe=boxfr,boxrule=0.5pt,arc=2pt,left=6pt,right=6pt,top=2pt,bottom=2pt}
\newcommand{\Cp}{\mathbb{C}}
\newcommand{\tr}{\operatorname{tr}}
\newcommand{\Cl}{C\ell_3}
\newcommand{\rv}[1]{\widetilde{#1}}
\newcommand{\cj}[1]{\overline{#1}}
\newcommand{\gr}[2]{\langle #1\rangle_{#2}}
\newcommand{\N}{\mathrm N}
\newcommand{\Zc}{\mathcal Z}
\newcommand{\Spin}{\mathrm{Spin}}
\newcommand{\rep}[1]{#1}
\newcommand{\forcesline}[1]{\par\vspace{-4pt}{\centering\small\textit{(#1)}\par}\vspace{2pt}}
\newcommand{\repline}[1]{\par\vspace{-2pt}{\leftskip=0.35in\rightskip=0.2in\noindent\textbf{Group class.} #1\par}\vspace{0pt}}
\newcommand{\lean}[1]{\texttt{\detokenize{#1}}}
\newcommand{\tg}[1]{{\scriptsize[#1]}}

\AtBeginDocument{\setlength{\abovedisplayskip}{2pt}\setlength{\belowdisplayskip}{2pt}\setlength{\abovedisplayshortskip}{2pt}\setlength{\belowdisplayshortskip}{3pt}}
\linespread{0.92}
\begin{document}
\color{black}\pdfrender{TextRenderingMode=FillStroke,LineWidth=0.22pt,StrokeColor=black,FillColor=black}
\vspace{0pt}
\begin{center}
{\Large\textbf{What in the $\mathbf{Hol}_{\boldsymbol\gamma}$ Bloch Ball, Batman?!}}\\[3pt]
{\large Or, Field Theoretic Self-Attention: A Canonicalization of Born, Lorentz, and Killing}\\[3pt]
{\large\textbf{Clifford edition}}\ \ {\normalsize written entirely in $Cl(3,0)$, the algebra of physical space}\\[3pt]
\textit{Companion to On Information Physics: From State Distinguishability to Causal Accumulation}
\end{center}
\vspace{1pt}
\begin{sheetbox}
\textbf{What this edition is.} The matrix edition restated in the real Clifford algebra $\Cl=Cl(3,0)$, generated by three vectors with $\sigma_j\sigma_k+\sigma_k\sigma_j=2\delta_{jk}$. Nothing is added or removed: $\Cl\cong M_2(\Cp)$ as real algebras, and every line below is the matrix line read through this dictionary. Matrices appear in one row and one diagram only.\par\vspace{3pt}
{\small\renewcommand{\arraystretch}{1.0}\setlength{\tabcolsep}{3pt}
\begin{minipage}[t]{0.5\linewidth}
\begin{tabular}[t]{@{}l l@{}}
\toprule
\textbf{Matrix edition} & \textbf{Clifford edition}\\
\midrule
$M_2(\Cp)$, 8 real dimensions & $\Cl$, grades 0--3 (dims 1, 3, 3, 1)\\
$i\,1$ & $I=\sigma_x\sigma_y\sigma_z$, central, $I^2=-1$\\
$\Cp\,1$ & centre $\Zc=\operatorname{span}_{\mathbb R}\{1,I\}$\\
$q^\dagger$ & reversion $\rv q$ (flips grades 2, 3)\\
$\operatorname{adj}q$ & Clifford conjugate $\cj q$ (flips grades 1, 2)\\
$\tfrac12\operatorname{tr}q$,\ \ $\tfrac12\operatorname{Re}\operatorname{tr}q$ & $\gr q0+\gr q3$,\ \ scalar part $\gr q0$\\
$\det q$ & norm $\N(q)=q\cj q\in\Zc$\\
\bottomrule
\end{tabular}
\end{minipage}\hfill
\begin{minipage}[t]{0.46\linewidth}
\begin{tabular}[t]{@{}l l@{}}
\toprule
\textbf{Matrix edition} & \textbf{Clifford edition}\\
\midrule
$\mathrm{Herm}(2)$ & paravectors, grades $0\oplus1$\\
$\mathfrak u(2)$;\ \ $\mathfrak{su}(2)$ & grades $3\oplus2$;\ \ bivectors\\
$GL(2,\Cp)$ & units $\Cl^\times$\\
$SL(2,\Cp)$ & $\Spin^+(1,3)=\{G:\N(G)=1\}$\\
$U(2)$ & unitary elements $\{U:U\rv U=1\}$\\
$SU(2)$ & rotors $\Spin(3)=\{R\in\Cl^+:R\rv R=1\}$\\
$U(1)$ & duality rotations $e^{I\varphi}$\\
\bottomrule
\end{tabular}
\end{minipage}}\par\vspace{3pt}
For $z=x+Iy\in\Zc$ write $\operatorname{Re}z=x$, $\operatorname{Im}z=y$, $|z|^2=z\rv z$. $\Cl^+$ is the even subalgebra (grades $0\oplus2$), the quaternions. \tg{sympy D1--D18}\par\vspace{2pt}
{\footnotesize\textit{Prior work.} $Cl(3,0)$ with paravectors as spacetime vectors is the algebra of physical space of W.~E.~Baylis (\textit{Electrodynamics: A Modern Geometric Approach}, 1999); the spacetime version is D.~Hestenes, \textit{Space-Time Algebra} (1966). The closure $J\rv J$ has the form of the current $\Psi\rv\Psi$ used there.}
\end{sheetbox}

\vspace{2pt}
\begin{sheetbox}
\textbf{Origin.} CAL is not a theory of quantum gravity. It is an information-physics engine for non-Markovian reality, built to lift self-attention to a field-theoretic form from first principles of cognition, causality and intelligence. Its derivations of known physics are unit tests: a law of causal memory that cannot recover established physics has no business modeling minds or societies.\par\vspace{2pt}
\textbf{Definition 1 (Information).} Information $:=$ distinguishability of physical states.\par\vspace{2pt}
\textbf{Remark 1 (Difference and observation).} To distinguish physical states is to observe a difference between them. Change is a difference between states along an ordered evolution; thus change is temporally ordered distinguishability. Observation here denotes physical distinguishability, not a required conscious observer.\par\vspace{2pt}
\textbf{Definition 2 (Causality).} Causality $:=$ change in physical state, including its rate of change.\par\vspace{2pt}
{\footnotesize\textit{From On Information Physics: From State Distinguishability to Causal Accumulation, Sec.~1.}}
\end{sheetbox}

\vspace{2pt}
\textbf{Motivation: nonlinear dynamics from linear kinematics.} The kinematics is linear: one carrier, with linear group actions. The Born rule, the Lorentz signature and the Killing form are obtained below kinematically, as quadratic readouts of the carrier and its closure, fixed before any equation of motion is written. The squaring is forced: no linear equivariant map carries the carrier to the paravectors, and $J\rv J$ is the only quadratic one, up to scale, invariant under right unitary elements \tg{sympy N1--N3}. CAL as printed is linear in $J$; its dynamics become nonlinear only where the closure re-enters the law as the context-dependent current, the \emph{self} in self-attention, and where it enters ($J$, $k$ or $\Pi$) is deferred. Admissible dynamics must respect the positive cone, the signature of the paravector slice and the invariant forms, so the kinematics is fixed first.\par
{\large\textbf{The Causal Accumulation Law}} (On Information Physics, Eq.~61)
\[
\begin{aligned}
T^{\mu\nu a}{}_{\rho\sigma}(x) \;=\; \underbrace{\alpha\, J^{\mu\nu a}{}_{\rho\sigma}(x)}_{\text{present source}} \;-\; (1-\alpha)\int_{J^-(x)} &\underbrace{k(\tau;x,x')}_{\text{history kernel}}\;\underbrace{\Pi^{\mu\nu}{}_{\rho\sigma|\alpha\beta}{}^{\gamma\delta}(x,x')}_{\text{slot coupling}}\\
\times\;\underbrace{P^{\alpha}{}_{\alpha'}(x,x')\,P^{\beta}{}_{\beta'}(x,x')\,P_{\gamma}{}^{\gamma'}(x,x')\,P_{\delta}{}^{\delta'}(x,x')}_{\text{spacetime transport of the past to }x}
&\;\times\;\underbrace{\Phi^{a}{}_{b}(x)\,P^{b}{}_{c'}(x,x')}_{\text{internal channel action}}\;\; J^{\alpha'\beta'c'}{}_{\gamma'\delta'}(x')\,d^4x'.
\end{aligned}
\]
{\large\textbf{Carrier.}}\quad $J=\sum_{a=0}^{3}Z_ae_a\in\Cl,\qquad Z_a=A_ae^{-H_a/\tau_E}+B_ae^{IS_a/\hbar_{\rm eff}}\in\Zc,\qquad A_a+B_a=1.$

\vspace{5pt}
{\large\textbf{Channels and multiplication}}\par\vspace{3pt}
\begin{minipage}[t]{0.6\linewidth}
\renewcommand{\arraystretch}{1.2}\setlength{\tabcolsep}{5pt}
\begin{tabular}[t]{@{}c c c c c c@{}}
\toprule
$a$ & $e_a$ & grade & $\rv{e_a}$ & paravector axis & under rotors\\
\midrule
$0$ & $1$ & 0 & $e_0$ & $1$ (time) & scalar $\mathbf1$\\
$1$ & $-I\sigma_x=\sigma_z\sigma_y$ & 2 & $-e_1$ & $Ie_1=\sigma_x$ & bivector $\mathbf3$\\
$2$ & $-I\sigma_y=\sigma_x\sigma_z$ & 2 & $-e_2$ & $Ie_2=\sigma_y$ & bivector $\mathbf3$\\
$3$ & $-I\sigma_z=\sigma_y\sigma_x$ & 2 & $-e_3$ & $Ie_3=\sigma_z$ & bivector $\mathbf3$\\
\bottomrule
\end{tabular}
\end{minipage}\hfill
\begin{minipage}[t]{0.36\linewidth}
\renewcommand{\arraystretch}{1.2}\setlength{\tabcolsep}{6pt}
\begin{tabular}[t]{@{}c|cccc@{}}
$\cdot$ & $1$ & $e_1$ & $e_2$ & $e_3$\\
\hline
$1$ & $1$ & $e_1$ & $e_2$ & $e_3$\\
$e_1$ & $e_1$ & $-1$ & $e_3$ & $-e_2$\\
$e_2$ & $e_2$ & $-e_3$ & $-1$ & $e_1$\\
$e_3$ & $e_3$ & $e_2$ & $-e_1$ & $-1$\\
\end{tabular}\par\vspace{2pt}
Row times column: $e_je_k=-\delta_{jk}+\epsilon_{jkl}e_l$.
\end{minipage}

\vspace{5pt}
The channel units are the scalar and the three unit bivectors; the coefficients $Z_a$ are scalar plus pseudoscalar, so $J$ fills all four grades. Each channel $a$ carries its own decay weight $A_ae^{-H_a/\tau_E}$ and phase weight $B_ae^{IS_a/\hbar_{\rm eff}}$; the phase is a duality rotation. The units are orthonormal, $\gr{e_a\rv{e_b}}0=\delta_{ab}$, so the channels add in quadrature and multiply as quaternions: $\gr{J\rv J}0=\sum_a|Z_a|^2$ is the Born weight, $\N(J)=J\cj J=\sum_aZ_a^2$ fixes purity (pure exactly when it vanishes), and the dual $Ie_j=\sigma_j$ turns the three bivector channels into the spatial axes of the paravector slice, with $e_0$ as time. \tg{Lean \lean{e1_e2}, \lean{trace_carrier_mul_conjTranspose}, \lean{det_carrier}; sympy Q1, B1, B2, D7, D8}

\vspace{3pt}\needspace{1.3in}
{\large\textbf{Born rule}}\quad $(\alpha=2,\ k=\delta(\tau-\tau_0),\ P\to1)$\hfill{\footnotesize\textit{Born rule paper, Thm.~6.1}}\par
\[
T\big|_{\alpha=2,\,\delta} = 2J(x)+J(x')
\]
\repline{\rep{Linear in $J$, so it stays in the carrier $\Cl=\Cl P_+\oplus\Cl P_-$, $P_\pm=\tfrac12(1\pm\sigma_z)$: two minimal left ideals, each a copy of the spinor $(\tfrac12,0)$ under left multiplication by $\Spin^+(1,3)$. With $P\to1$ the transport is the identity. \tg{Lean \lean{cal_alpha_two}; sympy B0, D16}}}

\[
\text{Born quadratic (closure): }\; J\rv J
\]
\forcesline{forces Born postulate; non-negativity from the signature}\repline{\rep{Invariant under right multiplication by unitary elements, $J\mapsto JU$; covariant under left $\Cl^\times$ by congruence $J\rv J\mapsto GJ\rv J\rv G$. The closure is reversion-even, so it is a paravector $t+\mathbf x$; its image is the closed future cone $t\ge|\mathbf x|$, interior $\Cl^\times/\{U\rv U=1\}$, null on the boundary. \tg{Lean \lean{born_right_unitary}, \lean{born_left_congruence}, \lean{born_posSemidef}; sympy B3--B7, D3, D9}}}

\[
\text{normalized: }\; \hat\varrho=\dfrac{J\rv J}{2\gr{J\rv J}0}=\tfrac12\big(1+\mathbf v\big)
\]
\forcesline{forces normalization}\repline{\rep{Quotient by dilations $\mathbb R_{>0}$; image is the Bloch ball $|\mathbf v|\le1$, and the Bloch vector $\mathbf v$ is literally a vector (grade 1). Rotors act on it as $SO(3)$ rotations, $\mathbf v\mapsto R\mathbf v\rv R$ (scalar $\mathbf1$ $\oplus$ vector $\mathbf3$); boosts do not preserve the scalar part. \tg{Lean \lean{su2_acts_by_rotation}, \lean{boost_changes_trace}; sympy L16, L17, D11}}}

\[
Z_a=A_ae^{-H_a/\tau_E}-B_ae^{IS_a/\hbar_{\rm eff}}\ \ \big(\theta=\tfrac{\alpha\pi}{2}=\pi\big)
\]
\repline{\rep{$e^{I\theta}$ is a duality rotation, a character of $U(1)=\{e^{I\varphi}\}$; at $\theta=\pi$ it is the sign rep of $\mathbb Z_2=\{\pm1\}$. \tg{Lean \lean{half_turn}, \lean{sign_character}}}}

\[
\gr{J\rv J}0=\sum_a\big|A_ae^{-H_a/\tau_E}-B_ae^{IS_a/\hbar_{\rm eff}}\big|^2
\]
\repline{\rep{$\gr{J\rv K}0$ is the Euclidean inner product on the eight real dimensions of $\Cl$, invariant under $J\mapsto UJ\rv V$ for unitary $U,V$. On the index $a$, two rotors $J\mapsto pJ\rv q$ act by real rotations, $(\Spin(3)\times\Spin(3))/\mathbb Z_2\cong SO(4)$, keeping both $\sum_a|Z_a|^2$ and $\N(J)=\sum_aZ_a^2$. \tg{Lean \lean{hs_invariant}, \lean{so4Equiv}, \lean{sumSq_left}; sympy B1, B8--B12, D8}}}

\[
p_m=\dfrac{\gr{E_mJ\rv J}0}{\sum_n\gr{E_nJ\rv J}0}=\dfrac{\gr{E_mJ\rv J}0}{\gr{J\rv J}0}
\]
\forcesline{forces additivity}\repline{\rep{Effects $E_m$ are paravectors in the closed future cone. Pairing of paravectors with their dual: invariant under $\varrho\mapsto G\varrho\rv G$, $E\mapsto\rv G^{-1}EG^{-1}$ (contragredient). $\sum_nE_n=1$ survives only for unitary $G$, so the $p_m$ are invariant under unitary elements. \tg{Lean \lean{born_kinematic}, \lean{born_additive}, \lean{povm_pairing}; sympy B13, B14, D5}}}

\vspace{4pt}

\vspace{3pt}\needspace{1.3in}
{\large\textbf{Lorentz signature}}\hfill{\footnotesize\textit{Born rule paper, Lemma~4.2}}\par
\[
\mathcal H(q)=\tfrac12\big(q+\rv q\big)=\gr q0+\gr q1
\]
\repline{\rep{Grade projector onto the paravectors in $\Cl=(\text{grades }0\oplus1)\oplus I\,(\text{grades }0\oplus1)$; commutes with $q\mapsto Gq\rv G$, so it is a $\Spin^+(1,3)$ intertwiner. \tg{Lean \lean{H_idem}, \lean{H_intertwines}; sympy L1--L3, D6}}}

\[
\Cl\xrightarrow{\;\mathcal H\;}\operatorname{span}_{\mathbb R}\{1,\sigma_x,\sigma_y,\sigma_z\}
\]
\repline{\rep{Under congruence $\Cl$ is two copies of the real $(\tfrac12,\tfrac12)$, grades $0\oplus1$ and grades $3\oplus2$; $\mathcal H$ keeps the first. \tg{Lean \lean{exists_herm}, \lean{A_eq_I_smul_H}; sympy L5, L6}}}

\[
X=t+\mathbf x,\qquad \mathbf x=x\,\sigma_x+y\,\sigma_y+z\,\sigma_z
\]
\repline{\rep{A paravector is a spacetime vector: real vector rep $(\tfrac12,\tfrac12)$, $X\mapsto GX\rv G$; kernel $\{\pm1\}$, so $SO^+(1,3)\cong\Spin^+(1,3)/\mathbb Z_2$. \tg{Lean \lean{spinor_kernel}; sympy L6, L9}}}

\[
\N(X)=X\cj X=(t+\mathbf x)(t-\mathbf x)=t^2-x^2-y^2-z^2
\]
\repline{\rep{$\N(GX\rv G)=|\N(G)|^2\,\N(X)$: invariant for $\Spin^+(1,3)$, conformal factor for $\Cl^\times$. Isometry group $SO^+(1,3)$. \tg{Lean \lean{det_herm}, \lean{det_congruence}; sympy L7, L8, D4, D9, D10}}}

\[
\sigma_j\longmapsto\text{Pauli matrices}:\qquad \Cl\cong M_2(\Cp)
\]
\repline{\rep{The one row with matrices: an isomorphism of real algebras, dimension 8, carrying $\rv q$ to $q^\dagger$, $\cj q$ to the adjugate, $\N$ to $\det$ and $2\gr q0$ to $\operatorname{Re}\tr$. In the quaternion notation of the matrix edition, $h=I$ and $h\mathbf i=\sigma_x$. \tg{Lean \lean{cliffordToM2_injective}, \lean{cliffordToM2_surjective}, \lean{paravector_norm}; sympy L10--L12, D2--D5}}}

\[
\hat\varrho=\tfrac12(1+\mathbf v):\quad 2\gr{\hat\varrho^{\,2}}0=1-2\,\N(\hat\varrho)=\tfrac12\big(1+v^2\big),\qquad \mathbf v=\mathbf x/t
\]
\repline{\rep{Purity, by squaring $\tfrac12(1+\mathbf v)$ with $\mathbf v^2=v^2$; $1-2\gr{\hat\varrho^{\,2}}0=1/(2\gamma^2)$. The Bloch radius is the velocity of the closure paravector: null ($v=1$) is pure, timelike ($v<1$) is mixed, the time axis ($v=0$, $\hat\varrho=\tfrac12$) is maximally mixed; the cone apex is zero. $\N$ is $\Spin^+(1,3)$-invariant but the scalar part $t$ is a time component, so purity is rotor-invariant and boost-dependent; only pure vs.\ mixed (null vs.\ timelike) is Lorentz-invariant. \tg{Lean \lean{purity_unit_trace}, \lean{purity_velocity}, \lean{purity_boost_dependent}; sympy L13--L15, D11}}}

\[
\text{Born weight}\equiv\text{signature selector: one grade projection }\mathcal H
\]
\repline{\rep{One intertwiner $\mathcal H$: the scalar part of its image is invariant under unitary elements (Born), the norm is $\Spin^+(1,3)$-invariant (Lorentz).}}

\vspace{3pt}\needspace{1.3in}
{\large\textbf{Purity from the channels}}\par
\[
v^2=1-\frac{|\N(J)|^2}{\gr{J\rv J}0^{\,2}}=1-\frac{\big|\sum_a Z_a^2\big|^2}{\big(\sum_a|Z_a|^2\big)^2}=\frac{4\sum_{a<b}\big[\operatorname{Im}(Z_a\rv{Z_b})\big]^2}{\big(\sum_a|Z_a|^2\big)^2}
\]
\repline{\rep{$\N$ is multiplicative and $\N(\rv J)=\rv{\N(J)}$, so $\N(J\rv J)=|\N(J)|^2$; then the purity row with $\gr{J\rv J}0=\sum_a|Z_a|^2$, $\N(J)=\sum_aZ_a^2$, and $|w|^2-\operatorname{Re}(w^2)=2(\operatorname{Im}w)^2$. Invariant under the two-rotor $SO(4)$ action on the channel index and a global duality rotation $Z_a\mapsto e^{I\varphi}Z_a$; the overall scale cancels. $v$ measures only the relative phase misalignment of the four channels. \tg{Lean \lean{purity_channels}, \lean{vsq_lagrange}, \lean{vsq_rotation}; sympy P1--P3, P6, D8}}}
\[
A_a=1:\quad Z_a=e^{-H_a/\tau_E}\in\mathbb R\ \Rightarrow\ J\in\Cl^+,\qquad J\rv J=J\cj J=\sum_aZ_a^2,\qquad \hat\varrho=\tfrac12,\qquad v=0
\]
\repline{\rep{Real $Z_a$ put $J$ in the even subalgebra, where reversion and Clifford conjugation coincide, so the closure equals the norm and is a scalar. The Boltzmann branch alone gives the maximally mixed closure, for every choice of $H_a$. \tg{Lean \lean{born_real_channels}, \lean{vsq_real}; sympy P4, D15}}}
\[
\tau_E\to0\ \ (H_a>0):\quad Z_a\to B_ae^{IS_a/\hbar_{\rm eff}},\qquad v=1\iff\sum_aB_a^2e^{2IS_a/\hbar_{\rm eff}}=0
\]
\repline{\rep{Only the phase branch can lift the closure to the light cone. Purity is a phase condition: $\N(J)=\sum_aZ_a^2=0$, that is, $J$ is null and lies in a single minimal left ideal. \tg{Lean \lean{boltzmann_limit}, \lean{vsq_eq_one_iff}; sympy P5}}}
\textbf{Scope.} These hold at fixed parameters. Scaling every $Z_a$ by one factor leaves $v$ unchanged; $v\to0$ when the real branch dominates every channel, or when a single channel dominates the sum. Adding the real branch does not lower $v$ monotonically in general, since a positive real offset can widen the angle between two channels. How $v$ evolves along a worldline needs equations of motion for $Z_a$, which are not yet written.\par

\vspace{3pt}\needspace{1.3in}
{\large\textbf{Generator slice (grades $3\oplus2$) and its Killing form}}\hfill{\footnotesize\textit{On Information Physics, Eq.~49}}\par
\[
\mathcal A(q)=\tfrac12\big(q-\rv q\big)=\gr q2+\gr q3
\]
\repline{\rep{Complementary grade projector, $\mathcal H+\mathcal A=1$; also commutes with $q\mapsto Gq\rv G$, so it is a $\Spin^+(1,3)$ intertwiner onto $I\,(\text{grades }0\oplus1)$. \tg{Lean \lean{A_antiHermitian}, \lean{A_intertwines}; sympy L2, L4, D6}}}

\[
\Cl\xrightarrow{\;\mathcal A\;}\operatorname{span}_{\mathbb R}\{I,\,e_1,\,e_2,\,e_3\}
\]
\repline{\rep{Pseudoscalar $\oplus$ bivectors $=\mathfrak u(1)\oplus\mathfrak{spin}(3)$, the reversion-odd elements, which generate the unitary elements. Under rotor conjugation this is $\mathbf1\oplus\mathbf3$ (the pseudoscalar is central, the bivectors turn as a vector). \tg{Lean \lean{exists_gen}, \lean{conj_splits_one_three}; sympy K1}}}

\[
A=y^0I+y^je_j,\qquad e_j=-I\sigma_j
\]
\repline{\rep{Generator coordinates; $e_j$ are the unit bivectors, $e_je_k=-\delta_{jk}+\epsilon_{jkl}e_l$. \tg{Lean \lean{e1_e2}; sympy Q1, D7}}}

\[
\N(A)=A\cj A=-(y^0)^2+|\vec y|^2
\]
\repline{\rep{Signature $(-,+,+,+)$. $A=IX$ with $X$ a paravector gives $\N(IX)=I^2\,\N(X)=-\N(X)$: duality (multiplication by the pseudoscalar) carries the state slice onto the generator slice, and the sign flip of the Minkowski form is $I^2=-1$. \tg{Lean \lean{det_gen}, \lean{det_I_smul}; sympy K2, K3, D12}}}

\[
\kappa(A,B)=\tr(\mathrm{ad}_A\mathrm{ad}_B)=8\big(\gr{A'B'}0+\gr{A'B'}3\big),\qquad A'=\gr A1+\gr A2
\]
\repline{\rep{Killing form of $\Cl$ under the commutator, as a complex Lie algebra ($\mathfrak{gl}(2)$); invariant under $\mathrm{Ad}$ (all Lie-algebra automorphisms). Only grades $1\oplus2$ enter: the centre drops out, so $\kappa$ is degenerate on the generator slice. \tg{Lean \lean{killingForm_gl2}, \lean{killing_ad_invariant}, \lean{killing_centre}; sympy K4, K5, K7, D13}}}

\[
\N(A)=|\vec y|^2=-A^2=-\tfrac18\,\kappa(A,A)\quad(y^0=0)
\]
\repline{\rep{On bivectors: $\kappa=8\gr{AB}0$ is negative definite (compact real form); the bivector algebra is simple, so by Schur the $\mathrm{Ad}$-invariant form is unique up to scale and the norm is that form. \tg{Lean \lean{det_su2_killing}, \lean{invariant_form_unique}; sympy K6, K8, D13}}}

\[
\text{grades }2\oplus1\ \text{under the commutator}\;\cong\;\mathfrak{so}(1,3)
\]
\repline{\rep{Real Killing form $16\gr{AB}0$ has signature $(3,3)$: rotations are the bivectors ($\kappa<0$) and boosts are the vectors ($\kappa>0$). The plane of vectors $u,v$ (paravectors $U,V$) is $A=\tfrac14(U\cj V-V\cj U)$; its area element is a Killing norm, $-\det h=\tfrac14\kappa(A,A)$, $h$ the Gram matrix. The norm is Schild's integrand; its root, real exactly on boost-like planes, is Nambu--Goto's, and buys reparametrization invariance. \tg{Lean \lean{killingForm_real}, \lean{killingForm_real_rotation}, \lean{killingForm_real_boost}; sympy K9, K10, D14, W1--W4, S1--S8}}}

\[
\N(X)\ \text{(Lorentz)}=\gr X0^2\ \text{(Born)}+\tfrac18\kappa(IX,IX)\ \text{(Killing)}\qquad(X\ \text{a paravector})
\]
\repline{\rep{Three readouts of one carrier: $\gr{\cdot}0$ is invariant under unitaries, $\N$ under $\Spin^+(1,3)$, $\kappa$ under $\mathrm{Ad}$. Invariant quadratics on the slice span two dimensions, so the constants are forced. Hence the vector part $\mathbf x$ has $|\mathbf x|^2=\gr X0^2-\N(X)=-\tfrac18\kappa(IX,IX)$: Killing reads the Bloch radius, equal to Born at rank one. A boost keeps $\N$ and trades $\gr{\cdot}0$ for $\kappa$, so the split is a choice of time axis. \tg{Lean \lean{threeway_su2}; sympy T1--T7}}}
\vspace{1pt}\needspace{2.4in}
{\large\textbf{The three readouts as one diagram}}\par\vspace{2pt}
\begin{minipage}[t]{0.495\linewidth}\centering
{\small\textbf{Matrix edition:\quad $c(J)=JJ^\dagger$}}\par\vspace{1pt}
$\begin{tikzcd}[column sep=3.6em,row sep=1.5em,every label/.append style={font=\small}]
M_2(\mathbb C) \arrow[r, "c"]
  & \mathrm{Herm}(2) \arrow[r, "{\operatorname{tr}}"] \arrow[dr, "{\det}"] \arrow[dd, "{i\,\cdot}"']
  & \mathbb R \ \ \text{(Born)}\\
  & & \mathbb R \ \ \text{(Lorentz)} \arrow[d, "-1"]\\
  & \mathfrak u(2) \arrow[r, "{\det}"']
  & \mathbb R \ \ \text{(Killing)}\\
  & \mathfrak{su}(2) \arrow[u, hook] \arrow[ur, "{-\frac18\kappa(A,A)}"']
\end{tikzcd}$
\end{minipage}\hfill
\begin{minipage}[t]{0.495\linewidth}\centering
{\small\textbf{Clifford edition:\quad $c(J)=J\widetilde J$}}\par\vspace{1pt}
$\begin{tikzcd}[column sep=3.6em,row sep=1.5em,every label/.append style={font=\small}]
C\ell_3 \arrow[r, "c"]
  & \text{grades }0\oplus1 \arrow[r, "{2\langle\cdot\rangle_0}"] \arrow[dr, "{\mathrm N}"] \arrow[dd, "{I\,\cdot}"']
  & \mathbb R \ \ \text{(Born)}\\
  & & \mathbb R \ \ \text{(Lorentz)} \arrow[d, "-1"]\\
  & \text{grades }3\oplus2 \arrow[r, "{\mathrm N}"']
  & \mathbb R \ \ \text{(Killing)}\\
  & \text{grade }2 \arrow[u, hook] \arrow[ur, "{-\frac18\kappa(A,A)}"']
\end{tikzcd}$
\end{minipage}\par\vspace{1pt}
Left and right are one diagram; every cell commutes. The closure lands in the paravectors, $\mathcal H\circ c=c$, so Born and Lorentz read one object. The square is $\N(IX)=-\N(X)$, the sign flip of the Minkowski form. The triangle is $\N(A)=-\tfrac18\kappa(A,A)$ on bivectors. So Lorentz and Killing are one norm $\N(q)=q\cj q$, restricted to grades $0\oplus1$ and to grade $2$; Born, a scalar part, is the only readout that depends on the observer. Because $c$, $\N$ and $\kappa(A,A)$ are not linear, the diagram lives among spaces with group actions, not in $\mathbf{Rep}(G)$: each space is a functor from $G$ (one object) to sets, each arrow a natural transformation, and $c(GJ)=G\,c(J)\,\rv G$ is the naturality square of the closure. \tg{Lean \lean{det_I_smul}, \lean{det_su2_killing}; sympy K3, K6, D12, D13}\par

\vspace{2pt}
\begin{smallbox}
\textbf{Remark (a coframe musing).} The purity row leaves a loose thread. The norm of the closure needs nothing but the carrier, yet its scalar part asks for a direction of time, and purity is the ratio of the two. On a curved background that direction would have to be supplied at every point, which is the job a coframe does in general relativity: the soldering $e^a{}_\mu\sigma_a$ turns $dx^\mu$ into the same paravectors used here, and its timelike leg $e_0$ plays the observer. Picture two observers sharing one event and one closure, one boosted relative to the other. They agree whether the closure sits on the light cone or inside it, so they agree whether it is pure; they disagree on how mixed it is, because each normalizes by their own time. Now let the closure be the memory term of CAL, $J(x')$ carried to $x$ by $P^\alpha{}_{\alpha'}$. The transport moves the closure, and a coframe at each end decides how much of it reads as certain. Whether that frame should be an extra input, or something the accumulated history itself selects, is open. It sits with the transport choices the Legend defers; a thought experiment only.
\end{smallbox}
\vspace{1pt}
\begin{smallbox}
\textbf{Remark (category theory, arrived at by accident).} The group-class notes were written as algebra and representation theory and turned out to be category theory without asking permission. Each step commutes with a group action. The linear steps (the grade projections $\mathcal H$, $\mathcal A$, and the scalar part for unitary elements) are morphisms of $\mathbf{Rep}(G)$, the intertwiners, and Schur's lemma is a statement about its $\mathrm{Hom}$-sets; the multilinear ones ($J\rv J$, $\N$, $\kappa$) are morphisms out of tensor products; the normalization $J\rv J/2\gr{J\rv J}0$ is nonlinear and is equivariant in the category of spaces with group actions. ``Canonical'' in the subtitle therefore means natural, in the technical sense: each readout is a natural transformation, and no basis or frame is chosen. The closure is the diagonal $J\mapsto(J,\rv J)$ then a bilinear product: its nonlinearity is all in the diagonal. Definition~1 lands in the same place. By the Yoneda lemma an object is determined up to isomorphism by $\mathrm{Hom}(-,X)$, how everything maps into it, and dually by $\mathrm{Hom}(X,-)$, how it maps out. States are the maps in, readouts the maps out, and the Born pairing composes the two; what distinguishes a state from every other is a map out. In short, you will be known by how you map, in and out. This is a reading of the definitions, not a theorem about CAL.\par\vspace{2pt}
\textbf{The invariance of the invariant.} This edition is that remark carried out. The carrier has three presentations, $\Cp\otimes\mathbb H$, $M_2(\Cp)$ and $\Cl$, isomorphic as real algebras, and the three readouts were defined without a basis, so each isomorphism carries the scalar part, the norm and $\kappa$ to themselves: the invariants are invariant under the change of algebra, a second naturality. By Yoneda this could not have failed, since isomorphic objects have the same maps; and Yoneda on Yoneda adds nothing, because the Yoneda embedding is fully faithful. The isomorphisms are checked; the rest is a reading. \tg{Lean \lean{carrier_coordsOf}, \lean{cliffordToM2_injective}, \lean{cliffordToM2_surjective}; sympy D2--D7}
\end{smallbox}

\vspace{1pt}
\begin{smallbox}
\textbf{Why Bell is violated}. CHSH: $S=E(a,b)-E(a,b')+E(a',b)+E(a',b')$. A fixed $\pm1$ answer sheet gives $S=A(B-B')+A'(B+B')=\pm2$, so $|S|\le2$. Two carriers give the joint closure $\varrho=\frac14\Big(1\otimes1+\sum_{j=1}^{3} e_j\otimes e_j\Big)$ in $\Cl\otimes_{\Zc}\Cl$, a sum over the three bivector channels. For unit vectors $\mathbf a$, $\mathbf b$, the scalar part of its product with $\mathbf a\otimes\mathbf b$ keeps the dot-product part: $E(a,b)=4\gr{(\mathbf a\otimes\mathbf b)\varrho}0=-\,\mathbf a\cdot\mathbf b=-\cos\theta_{ab}$.
The cosine, not the local line $-(1-2\theta/\pi)$, lets $|S|$ reach $|b+b'|+|b-b'|=2\sqrt2$ when $b\perp b'$. Each marginal alone is a fair coin. \tg{Lean \lean{chsh_local}, \lean{bell_correlation}, \lean{chsh_tsirelson}; sympy C1--C6, D17}
\end{smallbox}

\vspace{1pt}
\begin{smallbox}
\textbf{Legend}. $J\rv J$: Born quadratic, reciprocal at $\alpha=2$.\quad $J$: the observed present source current.\quad $\rv J$: the $J(x')$ of the earlier step, the memory (integrated past causal impulses) carried into the present, where the two interact. It works as a Bayesian update: $\rv J$ is the prior, $J$ the new observation, and the scalar part normalizes their product.\quad \textbf{Notation note:} the concept is settled; the notation is provisional. How the propagator and the gauge bundle enter is not yet fixed, so writing this pairing as $J\rv J$ may be in error and will be updated once that is worked out.\quad \textbf{Transport note:} the actual choices on how transport is handled are deferred until the author has properly learned it.\par\vspace{1pt}
\textbf{Uniqueness:} on the normalized image $J_\lambda=1+\lambda\mathbf n$ ($\lambda\in\Zc$, $|\lambda|=1$, $\mathbf n$ a unit vector, so $\rv{\mathbf n}=\mathbf n$, $\mathbf n^2=1$, $\gr{\mathbf n}0=0$), $J_\lambda\rv{J_\lambda}=2+(\lambda+\rv\lambda)\mathbf n$ and $\alpha J_\lambda=\alpha+\alpha\lambda\mathbf n$, so $\alpha J=J\rv J$ forces $\alpha=2$ and $\lambda=\pm1$, and $J/2=\tfrac12(1\pm\mathbf n)$ is a primitive idempotent. \tg{Lean \lean{uniq_alpha_dot}$\sigma$, \lean{uniq_projector}; sympy U1--U4} \textbf{Lemma disclosure:} the forward lemma is not yet written: that the $\alpha=2$ integration carries the present source current, through the normalization map, to $1-\mathbf n$. Until it is, $2J=J\rv J$ is a compatibility result, unique at $\alpha=2$, not a forward derivation from CAL; the identity admits both $\lambda=\pm1$, and only the half turn $e^{I\pi}=-1$ selects $\lambda=-1$.\par\vspace{1pt}
$A_ae^{-H_a/\tau_E}$, $B_ae^{IS_a/\hbar_{\rm eff}}$: decay and phase exponentials of $Z_a$.\quad $\theta=\alpha\pi/2$: fractional phase, and $e^{I\pi}=-1$ at $\alpha=2$, so $Z_a=A_ae^{-H_a/\tau_E}-B_ae^{IS_a/\hbar_{\rm eff}}$.\par\vspace{1pt}
$E_m$: detector effect, a paravector in the closed future cone, $\sum_mE_m=1$. The bottom of the Born rule is the total over all outcomes; $\sum_nE_n=1$ and the linear scalar part make it $\gr{J\rv J}0$.\quad $Z_a$: carrier component.\quad $\tau_E$: thermal width of the cost channel.\quad $\hbar_{\rm eff}$: phase scale.
\end{smallbox}

\vspace{1pt}
{\large\textbf{The necessity chain}}\hfill{\footnotesize\textit{each element is forced by a modeling requirement}}\par\vspace{2pt}
{\small\renewcommand{\arraystretch}{1.0}\setlength{\tabcolsep}{3pt}\setlength{\aboverulesep}{1pt}\setlength{\belowrulesep}{1.5pt}
\begin{longtable}{@{}>{\raggedright\arraybackslash}p{1.75in}>{\raggedright\arraybackslash}p{2.25in}>{\raggedright\arraybackslash}p{1.5in}>{\raggedright\arraybackslash}p{1.75in}@{}}
\rowcolor{hdr}\textbf{Requirement} & \textbf{Structure / Element} & \textbf{AI Analogy} & \textbf{Role}\\
\midrule\endfirsthead
\rowcolor{hdr}\textbf{Requirement} & \textbf{Structure / Element} & \textbf{AI Analogy} & \textbf{Role}\\
\midrule\endhead
Track state evolution & $J(x)$\ \ (source current) & Transformer-like & Context-dependent current/query\\
\multicolumn{4}{@{}>{\raggedright\arraybackslash}p{7.5in}@{}}{\hspace{0.25in}\textbf{Group class.} $J\in\Cl=\Zc\otimes\Cl^+$ (centre times quaternions); left action $J\mapsto GJ$ of $\Spin^+(1,3)$: two minimal left ideals, each $(\tfrac12,0)$; paravectors $X\mapsto GX\rv G$: vector $(\tfrac12,\tfrac12)$ of $SO^+(1,3)$ \tg{Lean \lean{carrier_coordsOf}, \lean{left_action_columns}; sympy L3, D16}}\\
\midrule
Compare states at different points & $\Pi^{\mu\nu}{}_{\rho\sigma\|\alpha\beta}{}^{\gamma\delta}(x,x')$\ \ (bitensor transport) & Geometric attention & Transport between states/positions\\
\multicolumn{4}{@{}>{\raggedright\arraybackslash}p{7.5in}@{}}{\hspace{0.25in}\textbf{Group class.} $\operatorname{End}(V_x)$, $V$ the rank-$(2,2)$ tensor rep of $SO^+(1,3)$, built from products of paravectors, $(\tfrac12,\tfrac12)^{\otimes2}=(0,0)\oplus(1,0)\oplus(0,1)\oplus(1,1)$; the antisymmetric part $(1,0)\oplus(0,1)$ is grades $1\oplus2$ \tg{Lean \lean{weights_vector_sq}, \lean{char_vector_sym}; sympy R1--R3}}\\
\midrule
Comparison must respect geometry & $P^{\alpha}{}_{\alpha'}(x,x')$\ \ (parallel propagators) & Relative positional encoding & Parallel transport in curved space\\
\multicolumn{4}{@{}>{\raggedright\arraybackslash}p{7.5in}@{}}{\hspace{0.25in}\textbf{Group class.} Holonomy in the frame-bundle structure group $SO^+(1,3)$ acting on paravectors; spinor lift in $\Spin^+(1,3)$ \tg{not machine-checked: deferred with the dynamics}}\\
\midrule
States must live in some space & $\Phi^a{}_b(x),\ P^b{}_{c'}(x,x')$\ \ (fiber holonomy) & Latent state geometry & Internal degrees of freedom\\
\multicolumn{4}{@{}>{\raggedright\arraybackslash}p{7.5in}@{}}{\hspace{0.25in}\textbf{Group class.} Internal gauge group $G_{\rm int}$ on the channel index $a=0,\dots,3$ (group not yet fixed); under rotor conjugation $a$ splits $\mathbf1\oplus\mathbf3$ ($e_0$ scalar, $e_j$ bivectors) \tg{Lean \lean{conj_splits_one_three}; sympy R5, R6}}\\
\midrule
History must fade appropriately & $k(\tau;x,x')$\ \ (memory kernel, proper time) & SSM-like & Memory kernel, history propagation\\
\multicolumn{4}{@{}>{\raggedright\arraybackslash}p{7.5in}@{}}{\hspace{0.25in}\textbf{Group class.} Lorentz scalar (depends on proper time $\tau$); convolution in $\tau$ is diagonalized by the characters $e^{s\tau}$ of $(\mathbb R,+)$, $s\in\Zc$, so $I^\alpha e^{s\tau}=s^{-\alpha}e^{s\tau}$ with $I^\alpha$ the fractional integral, not the pseudoscalar (lower limit $-\infty$, $\operatorname{Re}s>0$) \tg{Lean \lean{frac_integral_exp}; sympy R7}}\\
\midrule
Causal structure must be enforced & $J^-(x)$\ \ (past light cone integration) & Causal attention mask & Restrict to admissible past\\
\multicolumn{4}{@{}>{\raggedright\arraybackslash}p{7.5in}@{}}{\hspace{0.25in}\textbf{Group class.} The cone $\N(X)\ge0$, $\gr X0\le0$; preserved by the orthochronous group $SO^+(1,3)$; causal automorphisms of Minkowski space are Poincar\'e $\rtimes$ dilations (Zeeman) \tg{cited}}\\
\midrule
 & Integration over $J^-(x)$ & Recurrent / continuous memory & Incremental state update\\
\multicolumn{4}{@{}>{\raggedright\arraybackslash}p{7.5in}@{}}{\hspace{0.25in}\textbf{Group class.} Invariant measure $d^4x'$ on an invariant domain: the integral is an intertwiner, covariant integrand in, covariant field at $x$ out \tg{not machine-checked: belongs to the dynamics}}\\
\midrule
 & Symmetry and precompute & Structured models & Smaller vectors, exploit invariants\\
\multicolumn{4}{@{}>{\raggedright\arraybackslash}p{7.5in}@{}}{\hspace{0.25in}\textbf{Group class.} Isotypic decomposition and Schur's lemma: equivariant maps act blockwise on irreducibles, here grade by grade, cutting free parameters \tg{Lean \lean{invariant_form_unique}; sympy K8}}\\
\midrule
\end{longtable}}

\vspace{2pt}
{\footnotesize
\emergencystretch=3em \textbf{Machine verification.} Lean tags and sympy labels B--U check the matrix edition and carry over to $\Cl$ through the isomorphism and the dictionary, both Lean theorems (\lean{dict_N}); D1--D18 pass in sympy and a matrix-free suite reruns every sympy label, 104 of 104. Lean 4.35.0 with Mathlib v4.35.0: 294 theorems, no \texttt{sorry}, standard axioms only; PhysLean cross-checks the Pauli matrices and its own $Cl(3,0)$. Sympy, matrix edition: 67 of 67 exact; N1--N3, W1--W4, S1--S8, T1--T7 in sympy only (T1 also in Lean). Source: \texttt{\detokenize{github.com/JewJitsu11B/Liorhybrid}}.
\textbf{Not machine-checked:} Zeeman's theorem (cited); Clebsch--Gordan as a decomposition (characters and weights checked); that $\Spin^+(1,3)\to SO^+(1,3)$ is onto and that grades $2\oplus1$ form $\mathfrak{so}(1,3)$ (cited; kernel, Killing signature and the $SO(4)$ isomorphism checked); the complex Killing form in Clifford form (sympy only); minimal left ideals beyond D16; the dynamics (deferred); the forward lemma; Prior work, Scope and the remarks (readings).\par}

\end{document}
